\changecolours{ScoutGreen}
\section[Chave Parcial]{Chave Parcial (Discriminador) e Chaves Compostas}

% ------------------------------------------------------------------------------
\twocolframe[title={O Conceito de Chave Parcial (Discriminador)},leftcol=7.1cm,rightcol=6.9cm]{%
  \alert{O Que é a Chave Parcial?}\par
  \itemseps[topsepi=0.5mm,itemsepi=1mm,labelsep=2mm,leftmargini=3mm]
  \begin{itemize}
    \item \textbf{Definição:} Conjunto de atributos que distingue diferentes instâncias da entidade fraca \textbf{que estão associadas a uma mesma entidade proprietária}.
    \item \textbf{Unicidade Local vs. Global:}
      \begin{itemize}\setlength{\itemsep}{0.3mm}
        \item Localmente: É única apenas para aquele proprietário.
        \item Globalmente: Pode se repetir livremente sob proprietários distintos.
      \end{itemize}
    \item \textbf{Notação Gráfica:} Elipse conectada à entidade fraca com o nome do atributo \textbf{sublinhado por linha tracejada}.
  \end{itemize}
}{%
  \alert{Formação da Chave Composta Completa}\par
  \itemseps[topsepi=0.5mm,itemsepi=1mm,labelsep=2mm,leftmargini=3mm]
  \begin{itemize}
    \item \textbf{Equação da Identificação:}
      $$\text{PK}(\text{Fraca}) = \text{PK}(\text{Proprietária}) \cup \text{Discriminador}$$
    \item \textbf{Exemplo Corporativo:}
      \begin{itemize}\setlength{\itemsep}{0.3mm}
        \item Proprietária: \texttt{FUNCIONARIO}(\underline{\texttt{cpf}})
        \item Fraca: \texttt{DEPENDENTE}(\dashuline{\texttt{nome\_dep}})
        \item Chave Lógica: (\texttt{cpf}, \texttt{nome\_dep})
      \end{itemize}
    \item \textbf{Integridade Referencial:}
      No modelo relacional futuro, a PK da proprietária migra como Foreign Key (FK) compondo a Primary Key (PK).
  \end{itemize}
}

% ------------------------------------------------------------------------------
\begin{frame}{Exemplo Cotidiano: Redes de Hotéis e Quartos}
  \vspace{0.1cm}
  \begin{center}
    \tikz[scale=0.70, transform shape]{
      % Hotel (Entidade Forte)
      \node[draw=ScoutGreen, fill=ScoutGreen!15, thick, font=\bfseries\sffamily\scriptsize, minimum width=2.4cm, minimum height=0.65cm] (hotel) {HOTEL};
      \node[draw=ScoutNavy, fill=ScoutNavy!10, ellipse, font=\tiny\sffamily, above left=0.5cm and 0.1cm of hotel] (pk_hotel) {\underline{id\_hotel}};
      \node[draw=ScoutNavy, fill=ScoutNavy!10, ellipse, font=\tiny\sffamily, below left=0.4cm and 0.1cm of hotel] (nome_h) {nome\_hotel};
      \draw[thick, ScoutNavy] (hotel) -- (pk_hotel);
      \draw[thick, ScoutNavy] (hotel) -- (nome_h);

      % Relacionamento Identificador
      \node[diamond, draw=ScoutNavy, double, double distance=1.2pt, fill=ScoutNavy!15, thick, inner sep=2pt, font=\bfseries\sffamily\tiny, right=1.0cm of hotel] (rel) {POSSUI\_QUARTO};

      % Quarto (Entidade Fraca)
      \node[draw=ScoutGreen, double, double distance=1.2pt, fill=ScoutGreen!15, thick, font=\bfseries\sffamily\scriptsize, minimum width=2.4cm, minimum height=0.65cm, right=1.0cm of rel] (quarto) {QUARTO};
      \node[draw=ScoutNavy, fill=ScoutNavy!10, ellipse, font=\tiny\sffamily, above right=0.5cm and 0.1cm of quarto] (disc) {\dashuline{numero\_quarto}};
      \node[draw=ScoutNavy, fill=ScoutNavy!10, ellipse, font=\tiny\sffamily, below right=0.4cm and 0.1cm of quarto] (valor) {valor\_diaria};
      \draw[thick, ScoutNavy] (quarto) -- (disc);
      \draw[thick, ScoutNavy] (quarto) -- (valor);

      % Conexoes Chen
      \draw[thick, black!80] (hotel) -- node[above, font=\scriptsize\bfseries] {1} (rel);
      \draw[thick, double, double distance=1.5pt, black!80] (rel) -- node[above, font=\scriptsize\bfseries] {N} (quarto);
    }
  \end{center}
  \vspace{-0.1cm}
  \begin{table}[h]
    \centering
    \resizebox{0.78\textwidth}{!}{%
    \begin{tabular}{l l l l}
      \toprule
      \textbf{ID Hotel (Proprietário)} & \textbf{Nome do Hotel} & \textbf{Discriminador (Quarto)} & \textbf{Chave Composta Global} \\
      \midrule
      \texttt{H-01} & \texttt{Hotel Beira-Mar} & \texttt{101} & \texttt{(H-01, 101)} \\
      \texttt{H-01} & \texttt{Hotel Beira-Mar} & \texttt{102} & \texttt{(H-01, 102)} \\
      \texttt{H-02} & \texttt{Pousada do Sol} & \texttt{101} & \texttt{(H-02, 101)} \\
      \bottomrule
    \end{tabular}%
    }
  \end{table}
  \vspace{-0.1cm}
  {\centering\tiny \textbf{Observação:} O quarto 101 repete-se em hotéis distintos sem qualquer colisão de chave primária.\par}
\end{frame}
